% Research branch /root/kronecker, 21 November 2026.
% This fragment is designed for inclusion in the root manuscript.
\section{Triple rectangles: exact calibration and competing dyadic components}

\subsection{K-001: the classical three-dimensional calibration}
Bremner and Hu, \emph{Fundamental invariants for the action of
$SL_3(\mathbb C)^3$ on $3\times3\times3$ arrays}, arXiv:1112.2949v2,
Theorem~2.4, printed p.~5 (also Introduction, printed p.~1), attribute to Vinberg the polynomial-ring
identity
\[
 \mathbb C[(\mathbb C^3)^{\otimes3}]^{SL_3^3}
       =\mathbb C[I_6,I_9,I_{12}].
\]
Here the subscripts are ordinary polynomial degrees. Consequently the
regrading required in the question gives, exactly,
\[
 H^{\mathrm{ten}}_3(t)=\frac1{(1-t^2)(1-t^3)(1-t^4)}.                 \tag{K-001a}
\]
This is a cited classical result, not a novelty claim. In particular,
$\lim_{t\to i}(1-t^4)H^{\mathrm{ten}}_3(t)=(1-i)/4\ne0$.

For completeness, its twelve constituent polynomials are
$x^2/48+x/4+c_r$ for even $r$, and $x^2/48+x/8+c_r$ for odd $r$, where
\[
(c_0,\ldots,c_{11})=
\left(1,-\frac7{48},\frac5{12},\frac7{16},\frac23,-\frac7{48},
\frac34,\frac5{48},\frac23,\frac3{16},\frac5{12},\frac5{48}\right).
                                                               \tag{K-001b}
\]
One verifies this either by partial fractions in (K-001a), or by
multiplying the corresponding generating series by its denominator.
Every constituent has degree at most two and period dividing twelve;
therefore exact agreement at three arguments in each residue class
also proves the identity. The constituents at residues $1,5$ are
$(x-1)(x+7)/48$. They have negative constant coefficient and vanish at
$x=1$. After replacing $x$ by $x+1$, all twelve polynomials have
nonnegative coefficients, and each is strictly positive for integers
$x\ge1$. With GCT VI's Definitions 1.2.2 and 1.2.4 this proves
\[
                  p(k_3)=s(k_3)=1.                            \tag{K-001c}
\]
This direct calculation makes no assertion about polynomial bounds as
$n$ varies.

\subsection{K-002: four order-eight spectral types, with exact certificates}
Put $n=4$, $M=8$, and fix $\zeta=e^{2\pi i/8}$. A diagonal triple
with all eigenvalues in $\mu_8$ and determinant product $\zeta$ is
specified by three multisets $A,B,C$ of four residues. Its tensor
fixed space has support
\[
 S(A,B,C)=\{(i,j,k):A_i+B_j+C_k=0\pmod8\}.                       \tag{K-002a}
\]
Suppose its fixed space contains a nonzero polystable tensor. Its orbit cannot approach zero under the diagonal maximal torus of
$SL_4^3$, since that would also make it unstable for the full group.
The torus form of the Hilbert--Mumford criterion therefore puts zero
in the convex hull of its torus weights. Equivalently, the incidence
columns of its support admit nonnegative masses whose three
one-coordinate margins are all equal; rescale the common positive
margin to one. Extending these masses by zero gives such a solution
on (K-002a). This is a necessary condition, including when
eigenvalues repeat; it asserts no sufficiency in the repeated case.

Here is a complete \emph{finite-lift} calculation of that necessary
condition. Up to factor permutations, permutations within a factor,
and scalar shifts whose three exponents sum to zero, the feasible
triples are exactly the following four. Rows are ordered, and all
weights are normalized so their total sum is $1\pmod8$.
\[
\begin{array}{c|c|c|c}
 j&A&B&C\\\hline
1&(0,1,2,4)&(0,2,4,7)&(0,2,4,7)\\
2&(0,1,4,6)&(0,1,4,6)&(0,1,4,6)\\
3&(0,1,2,4)&(0,2,3,4)&(0,2,3,4)\\
4&(0,1,2,5)&(0,1,5,6)&(0,2,5,6)
\end{array}                                                     \tag{K-002b}
\]
All four spectra are simple. All four fixed spaces have ten tensor
coordinates; their incidence matrices have rank ten and Smith
nonzero diagonal $(1,1,1,1,1,1,1,1,1,8)$. All feasible mass vectors
are strictly positive and unique.

The exact certificate is \texttt{dyadic\_n4\_exact\_certificate.json};
the independent rational checker is
\texttt{verify\_n4\_certificate.py}. The exhaustive enumeration first
normalizes a supported triple of exponents to $(0,0,0)$, so every
multiset contains zero. There are $\binom{10}{3}=120$ such multisets.
It discards triples of even total exponent and those failing to
cover a coordinate. After the indicated symmetries there are exactly
2,285 cases. Four have the displayed positive rational solutions.
Each other case has an explicitly listed rational vector $y$ with
$A_S^ty\ge0$ and ${\bf1}^ty<0$ (with multiplicities in ${\bf1}$ when
weights repeat). Multiplying a proposed nonnegative solution by $y$
would contradict these two inequalities. The checker verifies all
2,281 separating vectors by rational arithmetic and separately
regenerates all 44,620 covered multiset triples. No floating-point
infeasibility assertion is used as a proof.

The scope of this statement is deliberate. It does not classify
all grading-fixed quotient points: a lift can have order larger than
eight while its determinant character has order eight. Such a lift
has not been excluded by this finite enumeration.

\subsection{K-003: a third smooth component, with positive local coefficient}
For the first spectral type in (K-002b), use the support tensor with
all ten nonzero coefficients equal to one. Its slices, symmetric in
the last two factors, are the quadratic space
\[
 Q=\langle q_0,q_1,q_2,q_3\rangle,
 \quad q_0=y_0^2+y_2^2,\quad q_1=2y_0y_3,\quad
 q_2=2y_1y_2+y_3^2,\quad q_3=2y_0y_2+y_1^2.                    \tag{K-003a}
\]
Its full product stabilizer consists of the two-dimensional
ineffective scalar torus times the diagonal $\mu_8$ given by the
first row of (K-002b). Its determinant character is faithful.

\paragraph{Proof of the full stabilizer.}
First the partial centroid of the four symmetric slice matrices is
scalar. If $EM_i=M_iE^t$ for all $i$, then $E$ preserves every
$\operatorname{im}M_i$. Intersecting the images of $M_0,M_1$ forces
$Ee_0=ae_0$. Symmetry with $M_0$ gives $Ee_2=de_2$.
Symmetry with $M_1$ gives $Ee_3=be_0+ae_3$; preservation of
$\operatorname{im}M_2=\langle e_1,e_2,e_3\rangle$ gives $b=0$.
Symmetry with $M_2$ gives $Ee_1=de_1+ce_2$.
Finally symmetry with $M_3$ gives $d=a$ and $c=0$.
Thus $E=aI$. Swapping the last two tensor factors now shows that any
product stabilizer has proportional last two matrices, reducing the
problem to the congruence stabilizer of $Q$.

The projective rank-at-most-two locus in $Q$ is
\[
 [q_0],\ [q_0+q_3],\ [q_0-q_3],\ [q_1].                       \tag{K-003b}
\]
For an exact elementary elimination check, take the sixteen $3\times3$
minors of the matrix of $\sum x_iq_i$. In the charts where $x_k=1$
and $x_{k+1}=\cdots=x_3=0$, their reduced lexicographic ideals are,
respectively, the zero ideal for $k=0$, $(x_0)$ for $k=1$, the unit
ideal for $k=2$, and $(x_0^2-1,x_1,x_2)$ for $k=3$.
These calculations are reproduced by \texttt{n4\_rank2.py}.
The first three points are collinear and the fourth is not. Hence a
congruence preserves $\mathbb Cq_1$ and $\langle q_0,q_3\rangle$.
The factor lines of $q_1$ are $\mathbb Cy_0,\mathbb Cy_3$; only the
first lies in the variable span $\langle y_0,y_1,y_2\rangle$ of
$\langle q_0,q_3\rangle$. Both lines are therefore individually fixed.
Write $y_0'=t_0y_0$, $y_3'=t_3y_3$.

Modulo $y_0$, the rank-one locus of the image of $Q$ consists precisely
of $\mathbb Cy_1^2,\mathbb Cy_2^2$: the other basis quadric is
$2y_1y_2+y_3^2$, and its matrix has rank at least two whenever its
coefficient is nonzero. Consequently $y_1',y_2'$ are scalar multiples
of a permutation of $y_1,y_2$, plus multiples of $y_0$.
Substitution in (K-003a) rules out the transposition, forces both
added multiples to vanish, and gives
\[
 t_2^2=t_0^2,\qquad t_1^2=t_0t_2,\qquad t_3^2=t_1t_2.
\]
After $t_0=1$, set $r=t_3$; then
$(t_1,t_2,t_3)=(r^6,r^4,r)$, $r^8=1$.
The first tensor factor is forced to be the inverse of the induced
slice action, up to scalar. Replacing $r$ by $r^{-1}$ gives exactly
the first spectral row of (K-002b). Their exponent sum is $33$, odd,
so the determinant character is faithful. This completes the proof.

The unique positive mass vector, ordered lexicographically by support
weights, is
\[
 (3/8,5/8,1/2,1/2,1/4,1/4,1/2,1/8,3/4,1/8).                  \tag{K-003c}
\]
Because the incidence matrix has rank ten, every choice of ten
nonzero coefficients is diagonally equivalent to every other over
$\mathbb C$. In particular the all-one tensor is equivalent to the
critical tensor with coefficients the square roots of (K-003c).
No two support triples coincide in two coordinates, so its three
reduced density matrices are identities. Kempf--Ness proves its
orbit closed. The faithful determinant character implies trivial
effective $SL_4^3$ stabilizer. Luna's slice theorem therefore makes
the quotient smooth there, of dimension $64-45=19$. The invariant
part of its tangent space has dimension $10-9=1$.
Thus its grading-orbit closure is a smooth-generic curve component
of the $\mu_8$ fixed locus, with normalized leading density one,
by the same numerical-semigroup argument as ALPOGE\_GCT P-F60.

\subsection{K-004: a smooth curve with the opposite local phase}
For the fourth spectral row, take the ten support triples
\[
\begin{split}
S={}&\{000,031,113,121,203,212,230,311,323,332\}.
\end{split}                                                    \tag{K-004a}
\]
Indices in this display are zero-based and refer to positions in
the sorted spectral rows, not to their numerical weights.
Its positive margin-one mass vector in this order is
\[
 (3/4,1/4,3/8,5/8,1/4,1/2,1/4,1/8,3/8,1/2).                 \tag{K-004b}
\]
Its full product stabilizer is again the ineffective scalar torus
times the diagonal $\mu_8$ described by the fourth row, with faithful
determinant character. Here is a proof that excludes nonmonomial
stabilizers.

Contract in the third factor, giving a $4\times4$ matrix $M(x)$.
The projective rank-at-most-two locus is
\[
 e_0,\quad e_2,\quad (r,r^2,0,1)\quad(r^4=1).                \tag{K-004c}
\]
The exact chart ideals of the $3\times3$ minors are the zero ideal,
the unit ideal, $(x_0,x_1)$, and
$(x_0^2-x_1,x_1^2-1,x_2)$ in charts $k=0,1,2,3$, respectively.
The plane $x_2=0$ is the unique plane containing five points of this
configuration; $e_2$ is the sixth point. Within that plane, $e_0$ is
the intersection of its two three-point lines. Any projective
transformation preserving (K-004c) therefore fixes $e_0,e_2$,
preserves this plane, and permutes the two pairs $r=\pm1$, $r=\pm i$.

Normalize its coefficient at $e_0$ to one. In the plane, the images
of $r$ and $-r$ have opposite first coordinate, so any additions of
$x_1,x_3$ to the transformed $x_0$ vanish. The possible $2\times2$
matrices on $(x_1,x_3)$ are the eight obtained from
\[
 (1,1)\longmapsto s^{-1}(s^2,1),\qquad
 (-1,1)\longmapsto it^{-1}(t^2,1),
 \quad s^4=t^4=1,\ t^2=-s^2.                                \tag{K-004d}
\]
The coefficient on $x_2$ is still an arbitrary nonzero $u$.

Now
\[
 \det M(x)=(x_1^2-x_3^2)(x_0^2-x_1x_3)-x_0x_1x_2^2.          \tag{K-004e}
\]
Preserving this quartic up to scalar eliminates the four
nondiagonal alternatives in (K-004d). For the other four,
$(x_1,x_3)$ has multipliers
$(1,1),(-1,-1),(i,-i),(-i,i)$, respectively, and
$u^2=1,-1,i,-i$, respectively. Thus the third-factor transformation,
up to scalar and inverse transpose, is exactly
\[
               \operatorname{diag}(1,r^2,r^5,r^6),\qquad r^8=1.
                                                               \tag{K-004f}
\]
Every such transformation occurs in the known diagonal stabilizer.
Divide by it. It remains to show that a stabilizer with third factor
$I$ is scalar. The four third slices are
\[
\begin{aligned}
M_0&=E_{00}+E_{23},&M_1&=E_{03}+E_{12}+E_{31},\\
M_2&=E_{21}+E_{33},&M_3&=E_{11}+E_{20}+E_{32}.
\end{aligned}                                                    \tag{K-004g}
\]
Their sum $M$ is invertible, with
\[
 M^{-1}=\begin{pmatrix}1&1&0&-1\\-1&0&1&0\\1&1&-1&0\\0&-1&0&1\end{pmatrix}.
\]
If $AM_iB^t=M_i$ for every $i$, then $A$ commutes with every
$M_iM^{-1}$. The simultaneous linear equations give
$A=aI$; an exact certificate is the 64-by-16 commutator system
in \texttt{n4\_stabilizer.py}, whose kernel is
$(a,0,0,0,0,a,0,0,0,0,a,0,0,0,0,a)$.
Alternatively multiplying the four explicitly displayed matrices
verifies the same equations directly. Hence $B=a^{-1}I$.
This proves the full stabilizer claim. Its weight sum is $33$, so
its determinant character is faithful. The balance, diagonal
normalization, free effective slice and fixed-curve arguments of
K-003 now apply verbatim, with (K-004b). In particular this is a
smooth-generic fixed curve, not merely a torus-semista-ble support.

\subsection{K-005: exact local residues and what their sum proves}
For any of the four simple spectra in (K-002b), let
\[
 C_j(\zeta)=
 \frac{\prod_{X=A,B,C}\prod_{a\ne b\in X}(1-\zeta^{a-b})}
 {\prod_{a\in A,b\in B,c\in C:\ a+b+c\ne0\ (8)}
       (1-\zeta^{a+b+c})}.                                   \tag{K-005a}
\]
The two new full-stabilizer proofs above, and ALPOGE\_GCT P-F60 and
P-F65 for rows 2 and 3, justify this as the actual moving conormal
factor of a smooth-generic fixed curve. With normalized density one,
ALPOGE\_GCT P-F23 supplies its local Lefschetz contribution.
Straight multiplication in $\mathbb Q[z]/(z^4+1)$ gives
\[
 (C_1,C_2,C_3,C_4)=
 \left(\frac1{256},\frac{1+\sqrt2}{128},
       \frac{3-2\sqrt2}{256},-\frac{\sqrt2}{256}\right).          \tag{K-005b}
\]
The computation is reproduced exactly by \texttt{dyadic\_local.py};
no numerical phase inference is used. In particular the fourth
component has the opposite principal-embedding phase from the
chain and P-F65 components. A proposed proof that all such
components have the chain phase is therefore false already at $n=4$.

Factor permutations produce respectively $3,1,3,6$ distinct curves.
Indeed their ordered spectra modulo scalar shifts distinguish them:
with a faithful determinant character and trivial effective
stabilizer, a prescribed grading generator has a unique lift modulo
the ineffective scalar kernel. An equivalence between two such
points would therefore identify their ordered spectra modulo those
scalar shifts. Enumerating those shifts and the six factor
permutations gives the stated orbit sizes. Thus there are at least
thirteen distinct smooth-generic curve components. Their combined
normalized coefficient is
\[
 3C_1+C_2+3C_3+6C_4=\frac{14-10\sqrt2}{256}\ne0.               \tag{K-005c}
\]
This is a sum over these thirteen components only. By itself this calculation is not a computation of the
residue of the complete Hilbert series, because further components
with higher-order lifts could contribute. K-007 and K-008 below
exclude that possibility in dimension four and upgrade (K-005c) to
the actual global residue.

\subsection{K-006: a support extension that doubles the character order}
Let $S$ be a support with independent incidence columns and positive
margin-one masses $m_e$. Suppose its diagonal stabilizer modulo the
ineffective scalar torus is cyclic of order $M$, and the determinant
character generates its character group. Select an edge
$e=(i,j,k)$ for which $Mm_e$ is odd, with $M$ a power of two.
Replace $e$ by
\[
 (i,j,N),\quad(i,N,k),\quad(N,j,k),\quad(N,N,N).                \tag{K-006a}
\]
The new masses on these four edges are
$m_e/2,m_e/2,m_e/2,1-m_e/2$. All other masses remain unchanged.
They are positive and have margin one. If the old support has no
pair of edges agreeing in two coordinates, the new support has the
same property. It therefore defines a critical, hence polystable,
tensor.

Here is the exact character-lattice proof of doubling. Let $s$ be the
character of the deleted edge. Eliminate the three new coordinate
characters using the first three new-edge relations. The fourth
relation becomes $2s=0$. Thus the new character lattice is obtained
from the old presentation by replacing the independent relation
$s=0$ by $2s=0$. Its torsion order is $2M$.
The determinant character changes by $-2s$, so it represents the
same element in this new group. The margin equations express
$M\chi=\sum_f Mm_f f$. Modulo the new relations this is $s$, since
$Mm_e$ is odd; this element is nonzero of order two. Hence $\chi$
has order exactly $2M$ and generates the entire torsion group.
The incidence columns remain independent by the same elimination.
This proves the claim for the full diagonal stabilizer, without a
claim about the full product stabilizer.

Apply this construction to K-004's mass-$3/8$ edge $(1,1,3)$,
then repeatedly to the newly added diagonal edge. After a scalar
gauge, the three lists of seed weights are
\[
 (-1,8,-7,-4),\qquad(-1,8,-4,5),\qquad(2,-4,-1,8),             \tag{K-006b}
\]
and all three lists acquire the same tail
\[
             -16,32,-64,\ldots,8(-2)^{n-4}\pmod{2^{n-1}}.
                                                               \tag{K-006c}
\]
The last weight is zero modulo $M=2^{n-1}$.
Their total exponent is $-7+16(-2)^{n-4}\equiv-7\pmod M$;
thus the determinant character is faithful of order $M$.
The eigenvalue-one tensor support consists exactly of the recursively
constructed $3n-2$ edges. To check this, the 64 seed sums are a fixed
finite calculation. Two tail weights and one seed weight cannot sum
to zero modulo $M$, by divisibility by sixteen. One tail weight and
two seed weights gives only $-16+8+8=0$, including its endpoint
interpretation when $M=16$. Three tail weights give precisely the
usual dyadic chain relations, by comparing their least two-adic
valuations. The small endpoint cases $n=4,5$ are checked directly.
Consequently the three spectra are simple and the fixed tensor
space has dimension $3n-2$ for this family.

This construction establishes actual polystable tensors and their
complete diagonal stabilizers. Full non-diagonal stabilizers, local
slice freeness, and global Hilbert residues for the extended family
remain separate questions. The files
\texttt{negative\_seed\_extensions.json} and
\texttt{extend\_support.py} record exact support and mass checks for
$n=4,\ldots,20$, and exploratory moving determinants. These numerical
determinants are not promoted to global pole claims.

\subsection{K-007: an a priori lift bound in dimension four}
The order-sixteen enumeration becomes a complete classification after
an elementary bound on the incidence lattice. Let
\[
 P_4=\{x\in\mathbb R_{\ge0}^{64}:\text{all twelve one-coordinate
 margins of }x\text{ equal }1\}.
\]
For a vertex $x$, let $S$ be its positive support, $m=|S|$, and
$B_S$ its $12\times m$ incidence matrix. Its columns are independent:
a dependence gives a small two-sided perturbation within the same
positive support, contradicting extremality. Hence $m\le10$.
The denominator $D(x)$ of all coordinates is exactly the order of
$\chi={\bf1}_{12}$ in
$\mathbb Z^{12}/B_S\mathbb Z^m$, since
$B_Sx=\chi$ and the columns are independent.

If some row has degree one, its unique edge has mass one; all other
edges meeting any of its three vertices have mass zero. It is an
isolated edge and deleting it reduces the problem to $P_3$.
Suppose there is no such row. Then $m\ge8$.
For $m=8$, all rows have degree two. The equations $x_e+x_f=1$
form a graph on the eight edge variables. An odd-cycle component
fixes all its variables to $1/2$; a bipartite component leaves one
free parameter. Since we have a vertex, there are no free components,
so $D(x)\le2$.

For $m=9$, every one of the three modes has three degree-two rows
and one degree-three row. Form the graph of degree-two equations.
Again each bipartite component supplies one parameter, while every
other component fixes its variables to $1/2$. On this affine solution
space the three remaining row sums are \emph{identical}: each is the
total mass of all nine variables minus three. They impose only one
additional equation. Extremality therefore permits at most one free
parameter. In that equation its coefficient is a nonzero integer
$a$ with $|a|\le3$, and the constant term is a half-integer. The
variable denominators consequently divide $2|a|$, so the two-primary
part of $D(x)$ is at most four.
The same argument handles the reduced $P_3$ problem: without a
further isolated edge its support size is six or seven. Six gives
only degree-two rows. Seven gives one degree-three row and two
degree-two rows in each mode, with the same single-equation argument.
Further isolated edges reduce to $P_2$, which is half-integral, or
$P_1$. Thus every vertex with $m<10$ has denominator with two-primary
part at most four.

For $m=10$, put $d=|\operatorname{tors}(
\mathbb Z^{12}/B_S\mathbb Z^{10})|$. The two row relations are the
block-sum differences. Their primitive two-by-two Pl\"ucker coordinates
are zero for two rows in the same mode and $\pm1$ otherwise.
Complementary-minor duality therefore says that the 48 nonzero
$10\times10$ minors of $B_S$ all have absolute value $d$.
Cauchy--Binet gives
\[
                   \det(B_S^tB_S)=48d^2.                     \tag{K-007a}
\]
The trace of $B_S^tB_S$ is thirty. Its largest eigenvalue is at least
$15/2$: evaluate its Rayleigh quotient on the all-one vector and use
$\sum_v\deg(v)^2\ge30^2/12$. The arithmetic--geometric mean inequality
for the other nine eigenvalues yields
\[
 48d^2\le \frac{15}{2}\left(\frac52\right)^9
       =\frac{15\,5^9}{2^{10}},\qquad d\le24.                 \tag{K-007b}
\]
The function $u((30-u)/9)^9$ decreases for $u\ge15/2$, justifying
use of the lower bound on the largest eigenvalue. In particular its
two-primary part is at most sixteen.

Now let a nonzero polystable tensor have a diagonal stabilizer with
determinant character of order divisible by eight. Its full support
admits a margin-one point, hence a vertex $x$ of that restricted
polytope. This is also a vertex of $P_4$, since the restricted polytope
is a face. The character order for the full support divides $D(x)$.
The preceding bounds force $m=10$. Adding the remaining support
relations makes the torsion group a quotient of that of $B_S$;
the free rank remains two, corresponding exactly to the ineffective
scalar torus. Its two-primary exponent is at most sixteen.
A diagonalizable complex algebraic group is a product of its torus
and a finite diagonalizable group, as follows by splitting its finitely
generated character group. The determinant character is trivial on
the torus. A prescribed primitive eighth-root value thus has a lift
of order at most sixteen in the finite two-primary factor.

This also applies to every nonvertex point of the grading fixed
scheme $X_4^{\mu_8}$. Choose its unique closed $SL_4^3$ orbit and a
product-linear lift of the grading generator. Equality of quotient
images supplies a stabilizer element with that determinant value.
Replace it by its semisimple part, which still stabilizes the tensor
and has the same character value, and diagonalize its three factors.
The preceding support argument replaces it by an order-at-most-sixteen
lift with the prescribed determinant character. No freeness of the
original stabilizer has been assumed in this reduction.

\subsection{K-008: the complete eighth-root fixed locus and pole for $n=4$}
Enumerate all multisets of four residues modulo sixteen that contain
zero. There are $\binom{18}{3}=816$. After coordinate and factor
permutations, supported scalar shifts, and normalization of the total
exponent to two, exactly 31,302 covered spectral orbits remain.
Of these, precisely four admit nonnegative margin-one masses.
Every feasible spectrum is even; after division by two and the
normalization of K-002, they are exactly its four rows. Each mass
solution is unique and strictly positive. For total exponent one
(character order sixteen), all 56,418 covered spectral orbits are
infeasible.

This is an exhaustive computation over integers, with rational
certificates. \texttt{enumerate16.cpp} generates every covered orbit
from the 816 multisets; it examines 1,229,228 covered triples for
character order eight and 2,333,632 for order sixteen before orbit
identification. The files
\texttt{m16\_chi8\_certificates.jsonl} and
\texttt{m16\_chi16\_certificates.jsonl} list four positive solutions
and 87,716 exact Farkas separating vectors. The independent checker
\texttt{verify16\_certificates.py} uses only rational arithmetic;
it verifies each certificate and matches every row of the spectral
lists. Thus no statement relies on the tolerance of the LP solver
that discovered the certificates.

By K-007 this enumeration covers every nonvertex point of
$X_4^{\mu_8}$. Strict positivity and uniqueness of the feasible mass
vectors force all ten support coefficients of any semistable tensor
in a listed fixed space to be nonzero. The rank-ten incidence matrix
makes the diagonal product group transitive on these nonzero
coefficient choices. Thus each spectral row gives exactly one
$GL_4^3$ orbit. The full-stabilizer and spectral arguments of
K-003--K-005 show that factor permutations give exactly thirteen
orbits, all pairwise distinct. The corresponding quotient points
are exactly thirteen nonzero grading orbits; their closures meet
only at the vertex. Every such curve is generically smooth in the
quotient and has normalized leading density one.
Consequently the entire positive-dimensional fixed locus has these
thirteen curve components. The zero-dimensional vertex and any
scheme structure supported there cannot affect the leading simple
pole coefficient.

ALPOGE\_GCT P-F23g explicitly gives the multiple-component version:
add its generic local terms over the top-dimensional fixed components.
Its Koszul/prime-filtration proof shows that all lower-dimensional
strata contribute lower pole order. Applying
it to the thirteen components proves the exact identity
\[
 \boxed{\displaystyle
 \lim_{t\to e^{2\pi i/8}}(1-t^8)H^{\mathrm{ten}}_4(t)
                  =\frac{14-10\sqrt2}{256}\ne0.}              \tag{K-008a}
\]
Thus the pole is simple and survives, with actual cancellation
between components of opposite phases. Galois conjugation proves
the corresponding statement at all four primitive eighth roots.

The exact part of the Hilbert series supported at the primitive
 eighth-root poles is
\[
             \frac{7-5t+5t^3}{256(1+t^4)}.                    \tag{K-008b}
\]
Indeed multiplying by $1-t^8$ and evaluating at a primitive eighth
root gives (K-008a) and its three conjugates, and these four simple
residues determine the strictly proper partial fraction part.
Its coefficient sequence is the constant-constituent periodic wave
\[
       \frac1{256}(7,-5,0,5,-7,5,0,-5),\quad\text{period }8.    \tag{K-008c}
\]
These are finite-dimensional calibration results. They do not,
by themselves, decide polynomial bounds on $p(k_n),s(k_n)$ or
numerical PH3 as $n$ varies. The complete rational Hilbert series
$H^{\mathrm{ten}}_4(t)$ has not been reconstructed here.
